\chapter{Travail principal}
\label{ch:main-work}

% Ce chapitre constitue le cœur du mémoire. Remplacez le petit algorithme
% ci-dessous par la construction, la méthode ou la suite de théorèmes
% principale. Énoncez les hypothèses avant les résultats et justifiez chaque
% étape.

\begin{assumption}[Sommets indexés]
\label{ass:indexed-vertices}
$V=\{1,\dots,n\}$, et les arêtes sont données par une matrice booléenne
$A\in\{0,1\}^{n\times n}$.
\end{assumption}

Sous l'\cref{ass:indexed-vertices}, la méthode classique de Warshall \cite{warshall1962}
met à jour une matrice booléenne de sorte qu'après l'itération $k$, elle décrive
les chemins dont les sommets intermédiaires appartiennent à
$\{1,\dots,k\}$. Une version compacte est donnée par
l'\cref{alg:warshall}.

\begin{algorithm}[htbp]
\caption{Clôture réflexive transitive d'une matrice d'adjacence}
\label{alg:warshall}
\begin{algorithmic}[1]
  \Require matrice d'adjacence $A\in\{0,1\}^{n\times n}$
  \Ensure matrice d'accessibilité $T$
  \State $T\gets A\lor I_n$
  \For{$k\gets 1$ à $n$}
    \For{$i\gets 1$ à $n$}
      \For{$j\gets 1$ à $n$}
        \State $T_{ij}\gets T_{ij}\lor(T_{ik}\land T_{kj})$
      \EndFor
    \EndFor
  \EndFor
  \State \Return $T$
\end{algorithmic}
\end{algorithm}

\begin{proposition}
\label{prop:warshall-correct}
À la fin de l'\cref{alg:warshall}, $T_{ij}=1$ exactement lorsque $j$ est
accessible depuis $i$.
\end{proposition}

\begin{proof}
On raisonne par récurrence sur $k$. Initialement, $T$ décrit les chemins de
longueur zéro ou un. À l'itération $k$, la mise à jour ajoute précisément les
chemins qui passent par le sommet $k$ ; après la dernière itération, $T$
représente donc la relation de l'\cref{eq:transitive-closure}.
\end{proof}
